\section{Scope, Method, and Qualifications}

\subsection{Reader, question, and claim boundary}

The intended reader is a serious general reader who can follow careful distinctions and is not presumed to know scholastic vocabulary, Latin, or the conventions of magisterial citation. Every technical term is glossed where it first does work.

The governing question is the relation between the particular and the general judgment: if the sentence passed at death is definitive, what is the general judgment for, and what may be claimed about either without going past the evidence. The governing claim, stated in the synthesis block in Section~\ref{sec:synthesis} and qualified there, is that the two are not two verdicts on one case but one verdict and its completion---differing in subject, in scope, and in publicity---and that the assessment of a man's desert is not reopened.

Theology governs the argument throughout. No canonical claim is made, no jurisdiction or body of law is in view, and no as-of date attaches to any discipline. The one disciplinary text quoted at length, Trent's decree on purgatory, is quoted for what it defines and for what it excludes from popular preaching, not as currently enforceable law.

\subsection{Series position}

This is the second of four independent Claude studies under the title \emph{The Four Last Things}: death, judgment, hell, heaven. The four were written concurrently and share no text. Where this volume must refer to another it does so by subject, never by page or section, because the companion volumes' pagination was not available to it. Death is presupposed here and not treated: what happens to a man at the moment of dying, and how Christians prepare for it, belong to the first volume. Hell is presupposed as the state into which the condemned sentence issues, and nothing is said here about its nature, its population, or what may be claimed about any person's presence in it. Heaven is presupposed as the state the blessed enter, and the vision of God is used only within the boundary \emph{Benedictus Deus} draws around it. Readers who want those subjects should go to their own volumes.

\subsection{Included and excluded material}

Included: the defined content of Lateran~IV, Lyons~II, \emph{Ne super his}, \emph{Benedictus Deus}, Florence, and Trent on the lot of the dead and on purgatory; the Catechism's treatment of the particular judgment, the resurrection, purgatory, and the Last Judgment; the 1979 letter of the Congregation for the Doctrine of the Faith on eschatology; the judgment scenes and sayings of Matthew, Luke, John, Paul, Hebrews, and the Apocalypse in the Douay--Rheims; Augustine's treatment of the last judgment in \emph{De civitate Dei} XX and of purgative fire in XXI; Aquinas's own treatment at \emph{Summa theologiae} III, q.~59; the Supplement's questions~69, 87, and 88 and its two appendices, with the provenance of that text established; and the corresponding place of Aquinas's \emph{Scriptum super Sententiis}, book~IV, distinction~47.

Excluded: the treatise on justification, merit, grace, and predestination, which is bounded in Section~\ref{sec:works} and not developed; Christology and soteriology, presupposed throughout; the nature of hell and of heaven, reserved to their own volumes; the intermediate state's philosophical problems (the mode of a separated soul's existence, whether temporal succession applies to it, personal identity across the interval); limbo and the fate of the unbaptized; the theology of indulgences and the whole question of how a suffrage is applied; private revelations about purgatory or the judgment, of which none is used or cited; the modern exegetical literature on the form and sources of Matthew 25:31--46, which was not surveyed; the Greek text of the New Testament, which was not examined at all; apocalyptic chronology, millennialism, and the identification of eschatological signs; and any judgment about the destiny of any named person.

\subsection{Evidence classes used}

\textbf{Defined teaching}: propositions asserted in the acts quoted in Section~\ref{sec:defined} and in Trent's canons, identified by document and Denzinger number at each appearance.

\textbf{Authoritative non-definitive teaching}: the Catechism's formulations, and the 1979 letter of the Congregation for the Doctrine of the Faith, whose status as an approved but non-definitive curial act is recorded where it is used.

\textbf{Scripture}: quoted only from the Douay--Rheims, in the tracked Challoner text of Project Gutenberg eBook~1581 registered in this repository's source library, with chapter and verse in the Vulgate numbering that text follows. Two quoted verses---Matthew 24:36 and John 5:22---stand in this repository's registered divergence table as differing from the settled American edition of 1899 that most printed Douay--Rheims Bibles reproduce; both divergences are recorded at their first appearance and neither affects an argument. No translation was composed for this article, and no other English version is quoted anywhere in it. Where the tracked text prints a typographic apostrophe it is typeset as an ordinary apostrophe.

\textbf{Patristic witness}: Augustine only, and only from \emph{De civitate Dei}. The English is Marcus Dods's 1871 printing in this repository's tracked transcription; three passages (XX.1, XX.14, XXI.26) were registered with this revision at exact artifact line bounds, and the complete chapters were read. One Latin check was made in the tracked TEI text of Hoffmann's CSEL~40.2. This is not a patristic survey and is not offered as one; the strength and the limits of that restriction are stated in the third overturning condition of the synthesis block.

\textbf{School theology}: Aquinas's own text at \emph{Summa theologiae} III, q.~59, distinguished throughout from the \emph{Supplementum Tertiae Partis}, whose status as a fourteenth-century transcription from Aquinas's commentary on the Sentences is established in Section~\ref{sec:trial} from the Leonine editors' own preface. Where the Supplement is used, that status is noted. Propositions taken from it are common or probable theological teaching, not defined doctrine, and several of them mark their own uncertainty in their own text; those markers are reproduced rather than suppressed.

\textbf{Project synthesis}: original to this article and labelled locally at each appearance as well as here. Namely: the observation that the fourfold enumeration has one slot where the acts have two events, and the use made of that; the observation that Lyons~II, \emph{Benedictus Deus}, and Florence all retain the general judgment in the same sentence in which they define immediate retribution, and the weight placed on \emph{nihilominus}; the formal reading of Matthew 25:31--46 as a non-parabolic prediction containing one simile, together with the use of Matthew~13 as a control case; the reading of Hebrews 9:27 as establishing the fixity of the life judged rather than the timing of the sentence; the adoption of the universal reading of \emph{omnes gentes}; the observation that ``particular judgment'' is a theological name for a defined immediacy rather than a separately defined event; the reading of the two surprised questions in Matthew~25 as incompatible with a ledger; the observation about the Catechism's non-repetition of the medieval Appendix's geography and identification of purgatorial with infernal fire; the verdict-and-vindication account of the two judgments in the synthesis block, with its stated concession; and the observation about the asymmetry between the Church's naming of the blessed and her silence about the damned. None of these is attributed to any source, and none should be repeated as the teaching of the Church.

\subsection{Method and its limits}

Every claim in this article rests on a source the authoring agent reached and read on 25~July 2026, with one recorded exception: the search that located Ecclesiasticus 7:40 in the tracked scriptural text was run on 26~July 2026 and its receipt is registered. Where a quotation is given, the wording was checked at the identified witness that day. Where a page image of a printed edition was inspected, that is stated; where only an electronic transcription or a machine transcription of a scan was available, that is stated too, and the weaker witness is never presented as the stronger.

Specific ceilings, all of which are recorded again in \texttt{research/source-audit.md}:

\begin{itemize}[leftmargin=1.2em,itemsep=1pt,topsep=2pt]
\item \textbf{Denzinger.} All conciliar and papal Latin except Trent was read in one dated web transcription of the Latin Enchiridion. That compilation is a finding aid to its underlying documents. No printed Denzinger edition was collated, and no conciliar manuscript or original register was consulted for any of Lateran~IV, Lyons~II, \emph{Ne super his}, \emph{Benedictus Deus}, or Florence. A reader who needs the exact promulgated wording of any of those acts should go behind this article to a critical text.
\item \textbf{Trent.} Both the English of Waterworth's 1848 translation and the Latin of the 1887 Tauchnitz stereotype reprint of the Roman edition were read at page images of the digitized copies, inspected directly rather than from optical character recognition. One orthographic variant between the Tauchnitz printing and the Denzinger text is recorded where it occurs. The Roman edition of 1834 that the Tauchnitz reprints was not itself consulted.
\item \textbf{Aquinas.} The English throughout is the Fathers of the English Dominican Province's, read in a dated web presentation; no printed copy of that translation was collated, and the English wording reported for \emph{Supplementum} q.~88, a.~1, ad~1 is that presentation's. The Latin of \emph{Summa theologiae} III, q.~59 and of \emph{Scriptum} IV, d.~47 was read in dated web transcriptions that identify the Leonine text but were not collated against it. The Latin of the Supplement was read partly at a Leonine page image (printed p.~205) and partly in the Internet Archive's uncorrected machine transcription of the Leonine volume, which is registered and used only as a finding aid; the comparison between the Supplement and the \emph{Scriptum} in Section~\ref{sec:trial} was made between two such witnesses and is not a collation.
\item \textbf{The Supplement's authorship.} Two incompatible accounts are on record and this article resolves neither. The Leonine preface dates the transcription to the middle of the fourteenth century and names only \emph{Compilator Supplementi}; a widely repeated attribution assigns it to Aquinas's socius Reginald of Piperno, who died about 1290. Both are reported. What is common to both, and all the argument requires, is that the Supplement is a compilation from the commentary on the Sentences and not a text Aquinas composed for the Summa.
\item \textbf{The Catechism.} Quoted from the Holy See's archived English IntraText presentation, whose delivery state was acquired and hashed on 25~July 2026. That presentation lowercases some sentence-initial articles and renders some typographic dashes as spaced hyphens; quotations here begin mid-sentence or restore standard casing only where the quoted words are unaffected. No printed edition was collated.
\item \textbf{The 1979 letter.} Quoted from the Holy See's English web text. The Latin published in \emph{Acta Apostolicae Sedis} 71 (1979) governs wording-critical claims and was not consulted.
\item \textbf{Greek.} Not examined. Nothing in this article turns on a Greek word, and the form-critical observations in Section~\ref{sec:matthew} are made about the Latin tradition's English and are labelled as this article's reading.
\end{itemize}

The article makes no negative-search claim about any corpus. Where it says that a phrase does not occur in the four acts of Section~\ref{sec:defined}, that is a statement about the wording of those acts as delivered by the registered witness, checked by reading them, not the result of a mechanical search over a bounded corpus.

\subsection{Rights}

The Douay--Rheims in Challoner's revision, Waterworth's 1848 English of Trent, the 1887 Tauchnitz Latin of Trent, Dods's 1871 English of the \emph{City of God}, Hoffmann's 1899--1900 CSEL Latin, the 1906 Leonine volume, and the English Dominican translation of the Summa are all public domain; the Latin acts of 1215, 1274, 1334, 1336, 1439, 1547, and 1563 are public domain by age. The Catechism and the 1979 letter are quoted briefly, with attribution, from official Holy See texts that remain outside this project's outbound licence; the exact bytes acquired from that host and from the two commercial transcription sites are recorded but not redistributed. The Denzinger compilation's numbering and apparatus are the editorial work of its successive editors. None of this material is relicensed by this publication, whose own prose and organization fall under the repository's project-content terms.

\subsection{Review state}

Completed: claim-driven research with same-day verification of every quoted witness; direct inspection of six printed page images (Trent English pp.~48, 49, 232, 233; Trent Latin pp.~38, 173) and two Leonine page images (the preface to the Supplement, and Supplement p.~205); registration of the new work, edition, artifact, and passage records listed in \texttt{research/source-bindings.toml}, with the human record in \texttt{research/scope.md} and \texttt{research/source-audit.md}; source-library validation; epistemic-order review, with received teaching labelled at every appearance and school opinion never presented as defined, audited in one place by the synopsis of claims by weight; boundary review against the other three volumes of the series; rights review; internal argumentative and quotation review; a multipass build with clean fatal-error, undefined-reference, and layout scans; the generation-metadata gate; the web-edition eligibility gate; and page-by-page raster review of the built PDF.

Outstanding: independent scriptural, patristic, Thomistic, dogmatic, theological, and ecclesiastical review. No imprimatur, nihil obstat, ecclesiastical approval, or completion of independent review is claimed, and internal checking is not independent review. The provenance finding in Section~\ref{sec:trial} in particular deserves the attention of someone who can read the Leonine apparatus properly and who has access to the manuscript scholarship this article did not consult.

Nothing here is pastoral counsel. A reader troubled about his own state or another's should take the matter to a confessor, not to a study document.
